\section{Praktikum}
\label{sec:praktikum-bab10}

\begin{praktikumbox}
\vspace{1.1em}
\textbf{Judul:} Menelusuri operasi aritmatika dan memeriksa
hasilnya.\\[0.3em]
\textbf{Sub-CPMK:} 3.1 \quad\textbullet\quad
\textbf{Alokasi waktu:} 100 menit \quad\textbullet\quad
\textbf{Bentuk:} individu
\end{praktikumbox}

\subsection{Tujuan Praktikum}

Setelah menyelesaikan praktikum ini, mahasiswa mampu menjalankan
instruksi aritmatika, mencatat perubahan register dan \eng{flag},
serta memverifikasi hasil operasi dengan DEBUG dan emu8086.

\subsection{Alat dan Bahan}

\begin{itemize}
    \item Komputer dengan emu8086 terpasang.
    \item DOSBox dengan TASM, TLINK, dan DEBUG.
    \item Berkas \file{operasi-aritmatika.asm} dari \file{code/asm/}.
    \item Lembar pengamatan register dan \eng{flag}.
\end{itemize}

\subsection{Langkah Kerja}

\begin{enumerate}
    \item \textbf{Merakit dan menjalankan program.}
    \begin{enumerate}
        \item Buka DOSBox, lalu masuk ke direktori kerja.
        \item Rakit dengan \file{tasm operasi-aritmatika.asm}.
        \item Link dengan \file{tlink operasi-aritmatika.obj}.
        \item Jalankan program. Program tidak mencetak keluaran.
    \end{enumerate}

    \item \textbf{Memeriksa hasil di segmen data.}
    \begin{enumerate}
        \item Jalankan \file{debug operasi-aritmatika.exe}.
        \item Ketik \file{d ds:0000 L 20} untuk melihat isi segmen
              data.
        \item Catat nilai \code{hasil\_lo}, \code{hasil\_hi},
              \code{kurang\_lo}, \code{kurang\_hi}, \code{kali},
              \code{bagi}, dan \code{sisa}.
        \item Bandingkan dengan Tabel~\ref{tab:hasil-program}.
    \end{enumerate}

    \item \textbf{Menelusuri dengan emu8086.}
    \begin{enumerate}
        \item Buka program pada emu8086.
        \item Jalankan \file{Single Step} sampai instruksi
              \lstasm|ADD AX, b_lo|.
        \item Catat \reg{AX}, \reg{BX}, dan \eng{flag} sebelum dan
              sesudah setiap instruksi.
        \item Ulangi untuk \lstasm|ADC AX, b_hi|,
              \lstasm|SUB AX, d_lo|, dan \lstasm|SBB AX, d_hi|.
    \end{enumerate}

    \item \textbf{Menguji MUL dan DIV.}
    \begin{enumerate}
        \item Catat \reg{AX} setelah \lstasm|MUL BL|.
        \item Catat \reg{AL} dan \reg{AH} setelah \lstasm|DIV BL|.
        \item Jelaskan perbedaan letak hasil bagi dan sisa.
    \end{enumerate}

    \item \textbf{Mengubah nilai untuk memicu flag.}
    \begin{enumerate}
        \item Ubah \code{a\_lo} menjadi \hex{7FFF} dan \code{b\_lo}
              menjadi \hex{0001}.
        \item Rakit ulang, lalu periksa \flag{CF}, \flag{OF}, dan
              \flag{SF} setelah \asm{ADD}.
        \item Ulangi dengan \code{a\_lo} = \hex{8000} dan
              \code{b\_lo} = \hex{8000}.
        \item Jelaskan mengapa \flag{CF} dan \flag{OF} dapat berbeda.
    \end{enumerate}

    \item \textbf{Menguji divide overflow.}
    \begin{enumerate}
        \item Ubah pembagi pada bagian \asm{DIV} menjadi 3 dan
              nilai yang dibagi tetap 1000.
        \item Hitung hasil bagi seharusnya.
        \item Jelaskan mengapa program harus memeriksa \reg{AH}
              sebelum \asm{DIV}.
        \item Tunjukkan baris mana yang mencegah kesalahan tersebut.
    \end{enumerate}
\end{enumerate}

\subsection{Hasil yang Dikumpulkan}

\begin{enumerate}
    \item Berkas \file{operasi-aritmatika.asm} versi modifikasi.
    \item Tabel pengamatan register dan \eng{flag}.
    \item Tangkapan layar isi segmen data pada DEBUG.
    \item Analisis perbedaan \flag{CF} dan \flag{OF}.
    \item Penjelasan syarat aman pembagian.
\end{enumerate}

\begin{asesmenbox}
\vspace{1.1em}
\textbf{Asesmen praktikum 8 (bobot 5\% dari komponen praktikum).}

\begin{tabular}{@{}p{6.2cm}cp{4.6cm}@{}}
    \toprule
    \textbf{Indikator} & \textbf{Bobot} & \textbf{Kriteria} \\
    \midrule
    Ketepatan hasil ADD/ADC dan SUB/SBB & 30\% &
    Nilai 32 bit dan flag tercatat benar \\
    \addlinespace
    Ketepatan hasil MUL dan DIV & 25\% &
    Hasil bagi, sisa, dan letak register benar \\
    \addlinespace
    Analisis flag & 25\% &
    Perbedaan CF dan OF dijelaskan dengan tepat \\
    \addlinespace
    Analisis divide overflow & 20\% &
    Syarat aman dan baris pencegahannya dijelaskan \\
    \bottomrule
\end{tabular}
\end{asesmenbox}
